\chapter{Carbon Dioxide (CO\textsubscript{2})}

\section*{What Is This Test?}
The blood carbon dioxide test measures the total carbon dioxide (CO$_2$) content in the serum or plasma. This value is often labeled total CO$_2$ (TCO$_2$) or bicarbonate (HCO$_3^-$) on chemistry panels. In blood, CO$_2$ exists in four forms: dissolved CO$_2$, carbonic acid (H$_2$CO$_3$), bicarbonate (HCO$_3^-$), and carbonate (CO$_3^{2-}$). The vast majority (approximately 90--95\%) is in the form of bicarbonate, which is why the total CO$_2$ value closely approximates the serum bicarbonate concentration. Total CO$_2$ is a measure of the metabolic component of acid-base balance. It is part of the basic and comprehensive metabolic panels (BMP/CMP). Carbon dioxide is produced by cellular metabolism, transported in the blood as bicarbonate, and eliminated by the lungs (as CO$_2$) and kidneys (as bicarbonate reabsorption and hydrogen ion excretion). The measurement of total CO$_2$ (bicarbonate) is used with the blood pH and pCO$_2$ (from arterial or venous blood gas analysis) to evaluate metabolic acidosis, metabolic alkalosis, and respiratory acid-base disorders, and to calculate the anion gap.

\section*{Alternative Names}
\begin{itemize}
  \item Total CO$_2$
  \item TCO$_2$
  \item Bicarbonate (HCO$_3^-$)
  \item Carbon Dioxide, Serum
  \item CO$_2$ (Blood)
  \item Venous CO$_2$
\end{itemize}
\section*{Principle}
Total CO$_2$ is measured on automated chemistry analyzers by an enzymatic method or an ion-selective electrode (ISE). The enzymatic method uses phosphoenolpyruvate carboxylase (PEPC), which converts phosphoenolpyruvate and CO$_2$ to oxaloacetate; the oxaloacetate is then measured by a coupled enzyme reaction (MDH/NADH), and the change in NADH absorbance is proportional to the CO$_2$ concentration. Alternatively, ISE methods use a CO$_2$-permeable membrane and measure the pCO$_2$ with a Severinghaus electrode. The value is reported in mmol/L (mEq/L). The total CO$_2$ is the sum of dissolved CO$_2$, carbonic acid, bicarbonate, and carbonate: TCO$_2$ = [HCO$_3^-$] + [CO$_3^{2-}$] + [H$_2$CO$_3$] + [dissolved CO$_2$]. In plasma, [HCO$_3^-$] $\approx$ TCO$_2$ $-$ 1.2 mmol/L. The relationship between CO$_2$, bicarbonate, and pH is described by the Henderson-Hasselbalch equation: pH = 6.1 + log ([HCO$_3^-$] / (0.03 $\times$ pCO$_2$)). This equation underlies the interpretation of acid-base disorders. Bicarbonate is the major buffer of the blood, and its concentration is regulated by the kidneys: the kidneys reabsorb virtually all filtered bicarbonate and generate new bicarbonate by excreting hydrogen ions (as titratable acid and ammonium, NH$_4^+$). The lungs regulate CO$_2$ through ventilation.

\section*{History}
The relationship between carbon dioxide, pH, and ventilation was established in the early 20th century. Lawrence J. Henderson derived the equilibrium equation for CO$_2$ and bicarbonate in 1908, and Karl Hasselbalch reformulated it into the logarithmic form (the Henderson-Hasselbalch equation) in 1916. The concept of the anion gap was introduced in the 1970s to categorize metabolic acidosis. Modern automated analyzers measure total CO$_2$ by enzymatic or ISE methods as part of the basic metabolic panel. The classification of acid-base disorders (respiratory vs. metabolic; simple vs. mixed) was formalized with the work of Robert Narins and others in the 1980s--1990s \cite{seifter2014integration}. Clinical guidelines for the evaluation of metabolic acidosis, including the measurement of serum osmolality, the osmolal gap, the anion gap, and the delta-delta, continue to guide practice \cite{kraut2007serum}.

\section*{Why This Test Is Ordered}
\begin{itemize}
  \item Evaluation of metabolic acidosis: total CO$_2$ (bicarbonate) is measured in patients with suspected metabolic acidosis (diabetic ketoacidosis, lactic acidosis, renal failure, diarrhea, salicylate poisoning, toxins such as methanol and ethylene glycol). A low bicarbonate (<22 mmol/L) with a widened anion gap or normal anion gap guides diagnosis
  \item Evaluation of metabolic alkalosis: a high bicarbonate (>29 mmol/L) is seen in metabolic alkalosis (vomiting, diuretics, hyperaldosteronism). Confirmation requires blood gas analysis and measurement of urine chloride
  \item Monitoring of chronic kidney disease and renal tubular acidosis: the bicarbonate level is part of the metabolic panel; a low bicarbonate indicates metabolic acidosis that may require alkali therapy (target serum bicarbonate >22 mmol/L)
  \item Assessment of the acid-base status in critically ill patients: total CO$_2$ from a venous sample provides an estimate of bicarbonate, which is used with the blood gas to determine the primary disorder and compensation
  \item Routine health screening: total CO$_2$ is a component of the basic and comprehensive metabolic panels (BMP/CMP), which are ordered for general health assessment, medication monitoring, and pre-operative evaluation
\end{itemize}

\section*{Primary vs Secondary Test}
Total CO$_2$ (bicarbonate) is a primary component of the metabolic panel and an essential part of acid-base assessment. It is used as a secondary test when a full blood gas (arterial or venous) with pH and pCO$_2$ is needed to characterize the acid-base disorder precisely.

\section*{How This Test Is Performed}
A healthcare professional draws blood from a vein into a serum separator tube (gold-top) or a lithium heparin tube (green-top). The sample should be transported to the laboratory promptly, because CO$_2$ is lost from the sample if the tube is left open or processed slowly (the sample must be capped and analyzed within 2 hours; ideally 30--60 minutes). The sample is centrifuged, and the serum or plasma is analyzed by the enzymatic or ISE method. Results are typically available within 1--2 hours. Arterial blood gas analysis measures the true pCO$_2$ and pH directly; the serum total CO$_2$ from the chemistry panel is used for the bicarbonate estimate.

\section*{What Preparation Is Needed}
No fasting is required. The patient should avoid hyperventilation (which lowers CO$_2$) immediately before the blood draw. The sample must be collected anaerobically (filled tube, not shaken) and processed promptly to prevent loss of CO$_2$ to the atmosphere. Document all medications, especially diuretics, bicarbonate therapy, corticosteroids, and carbonic anhydrase inhibitors (acetazolamide).

\section*{Contraindications and Precautions}
\begin{itemize}
  \item Total CO$_2$ is a bicarbonate estimate, not a direct measurement of blood pH or pCO$_2$. A normal total CO$_2$ does not exclude a mixed acid-base disorder (e.g., coexisting metabolic acidosis and metabolic alkalosis may normalize the bicarbonate)
  \item Sample handling errors: exposure of the sample to air (uncapped tube) causes loss of CO$_2$ and falsely low total CO$_2$. Delayed centrifugation or an open tube falsely lowers the value. Hemolysis can interfere with some enzymatic assays
  \item The anion gap is calculated as Na$^+$ - (Cl$^-$ + HCO$_3^-$). A low total CO$_2$ (low bicarbonate) widens the anion gap; interpretation requires sodium, chloride, and glucose measurements from the same sample
  \item Total CO$_2$ does not reflect the respiratory component of acid-base balance. A patient with a normal bicarbonate may have a severe respiratory acidosis or alkalosis (abnormal pCO$_2$); blood gas analysis is required to detect these
  \item Bicarbonate is affected by potassium: hyperkalemia can cause a mild metabolic acidosis, and correction of the potassium may correct the acidosis. Serum potassium should be interpreted with the bicarbonate
  \item Drugs alter bicarbonate: acetazolamide (carbonic anhydrase inhibitor) causes metabolic acidosis; sodium bicarbonate, diuretics, and corticosteroids can cause metabolic alkalosis. A patient receiving IV sodium bicarbonate will have a high total CO$_2$ regardless of the underlying disease
\end{itemize}
\section*{Risks}
Standard venipuncture risks: mild pain, bruising, or hematoma at the puncture site.

\section*{What Happens After the Test}
The total CO$_2$ is interpreted with the serum sodium, chloride, glucose, BUN, creatinine, and the blood pH and pCO$_2$ when a blood gas is available. The anion gap is calculated: AG = Na$^+$ - (Cl$^-$ + HCO$_3^-$), normal 8--12 mmol/L (with albumin correction). A widened anion gap acidosis (methanol, uremia, DKA, propylene glycol, iron/isoniazid, lactic acidosis, ethylene glycol, salicylates) requires specific diagnosis and treatment. A normal anion gap acidosis (diarrhea, renal tubular acidosis, acetazolamide, ureterosigmoidostomy) requires measurement of urine anion gap and urinary pH. A high bicarbonate suggests metabolic alkalosis and prompts evaluation of urine chloride.

\section*{Understanding Results}

\subsection*{Below Normal}
\begin{itemize}
  \item \textbf{Low total CO$_2$ (<22 mmol/L) with a widened anion gap (>12 mmol/L):} High anion gap metabolic acidosis. Causes (the mnemonic MUDPILES or GOLDMARK): methanol, uremia, DKA, propylene glycol, iron/isoniazid, lactic acidosis, ethylene glycol, salicylates. Lactate is measured when lactic acidosis is suspected; glucose and ketones for DKA; osmolal gap and toxicology screen for methanol/ethylene glycol
  \item \textbf{Low total CO$_2$ with a normal anion gap (8--12 mmol/L):} Normal (hyperchloremic) anion gap metabolic acidosis. Causes: diarrhea, renal tubular acidosis (distal and proximal), acetazolamide, ureterosigmoidostomy, early renal failure, HCl or ammonium chloride ingestion. The urinary anion gap and urinary pH help distinguish the subtypes
  \item \textbf{Low total CO$_2$ with respiratory alkalosis:} Chronic hyperventilation (hypoxia, high altitude, anxiety, sepsis, salicylate poisoning early phase) lowers pCO$_2$, and the kidney compensates by reducing bicarbonate reabsorption. The anion gap may help identify salicylate poisoning (early mixed picture)
  \item \textbf{Low total CO$_2$ in chronic kidney disease:} Bicarbonate falls as the kidney loses the ability to excrete acid (eGFR <30 mL/min). This metabolic acidosis contributes to bone disease and muscle wasting; alkali therapy (target bicarbonate >22 mmol/L) is recommended
\end{itemize}

\subsection*{Above Normal}
\begin{itemize}
  \item \textbf{High total CO$_2$ (>29 mmol/L) with hypochloremia:} Metabolic alkalosis with chloride depletion. The most common causes are vomiting (gastric acid loss) and loop or thiazide diuretics. The urine chloride is low (<10--20 mmol/L) in chloride-responsive alkalosis (vomiting, diuretics, post-hypercapnia); a high urine chloride (>20 mmol/L) suggests chloride-resistant alkalosis (hyperaldosteronism, Cushing syndrome, licorice, Bartter/Gitelman syndromes)
  \item \textbf{High total CO$_2$ with normal chloride:} Consider alkali administration (bicarbonate, citrate, antacids), or mineralocorticoid excess with a normal chloride. The cause is usually evident from the history and medications
  \item \textbf{High total CO$_2$ with chronic respiratory acidosis:} Chronic CO$_2$ retention (COPD, obesity hypoventilation, chest wall disease) stimulates renal bicarbonate retention; the bicarbonate rises to compensate. Compensation is incomplete, so the pH is slightly acidotic; pCO$_2$ is elevated on blood gas
  \item \textbf{High total CO$_2$ with high sodium:} Volume depletion with contraction alkalosis (diuretics, vomiting) or hyperaldosteronism (Conn syndrome). Sodium, potassium, chloride, and urine chloride help distinguish the cause
\end{itemize}

\section*{Normal Results}
\begin{itemize}
  \item \textbf{Total CO$_2$ (adults):} 22--29 mmol/L (mEq/L) in venous serum/plasma
  \item \textbf{Bicarbonate (HCO$_3^-$):} approximately 1--1.2 mmol/L lower than total CO$_2$ (the difference is dissolved CO$_2$); reported interchangeably in most panels
  \item \textbf{Arterial pCO$_2$:} 35--45 mmHg (normal reference for the respiratory component, measured on blood gas)
  \item \textbf{Arterial pH:} 7.35--7.45 (measured on blood gas; not derived from total CO$_2$)
\end{itemize}

\section*{Follow-up Tests}
\begin{itemize}
  \item \textbf{Arterial (or venous) blood gas with pH and pCO$_2$:} Required to characterize the acid-base disorder. The Henderson-Hasselbalch equation links pH, pCO$_2$, and bicarbonate. Compensation rules (Winter's formula for metabolic acidosis: expected pCO$_2$ = 1.5 $\times$ [HCO$_3^-$] + 8 $\pm$ 2) identify mixed disorders
  \item \textbf{Anion gap and delta-delta (delta AG / delta HCO$_3^-$):} Calculated from sodium, chloride, and total CO$_2$. A delta-delta <1 indicates a coexisting normal anion gap acidosis; >2 indicates a coexisting metabolic alkalosis
  \item \textbf{Serum lactate:} To diagnose lactic acidosis in a patient with a widened anion gap acidosis. Normal <2 mmol/L; levels >4 mmol/L with acidosis are consistent with type A or type B lactic acidosis
  \item \textbf{Serum glucose, ketones, BUN, creatinine, and osmolality:} To distinguish diabetic ketoacidosis (glucose and ketones), uremic acidosis (BUN/creatinine), and methanol/ethylene glycol poisoning (osmolal gap >10 mOsm/kg)
  \item \textbf{Urine chloride and urine anion gap:} Urine chloride distinguishes chloride-responsive from chloride-resistant metabolic alkalosis. The urine anion gap (urine Na + urine K - urine Cl) distinguishes GI losses (positive gap) from renal tubular acidosis (negative gap)
  \item \textbf{Serum potassium and magnesium:} Hypokalemia frequently accompanies metabolic alkalosis and requires correction; hyperkalemia causes a mild metabolic acidosis. Electrolyte correction is often needed before the acid-base disorder will correct
\end{itemize}

\section*{References}
\bibliographystyle{unsrtnat}
\bibliography{references}

\clearpage
\section*{Regulatory Status and Authorization Date}
This chapter describes a test category, not a specific commercial product. FDA, CDSCO, EU IVDR, and MHRA approval or registration dates therefore cannot be assigned without the assay manufacturer, product name, instrument, and intended use. For laboratory-developed tests, an individual agency approval date may not exist. See the Preface for the interpretation of regulatory status.

\clearpage
